% Generated draft: requires mathematical and source review.
\documentclass[11pt]{article}
\usepackage[margin=25mm]{geometry}
\usepackage{fontspec}
\setmainfont{Latin Modern Roman}
\usepackage{amsmath,amssymb,mathtools,unicode-math}
\setmathfont{Latin Modern Math}
\usepackage{enumitem,needspace,xurl,hyperref}
\setlength{\parindent}{0pt}
\setlength{\parskip}{5pt}
\setlength{\emergencystretch}{3em}
\newcommand{\N}{\mathbb{N}}
\newcommand{\Z}{\mathbb{Z}}
\newcommand{\Q}{\mathbb{Q}}
\newcommand{\R}{\mathbb{R}}
\newcommand{\C}{\mathbb{C}}
\newcommand{\F}{\mathbb{F}}
\newcommand{\A}{\mathbb{A}}
\newcommand{\PP}{\mathbb{P}}
\newcommand{\E}{\mathbb{E}}
\newcommand{\eps}{\varepsilon}
\newcommand{\dd}{\,\mathrm{d}}
\begin{document}
\section*{UnsolvedMath ID: 30003866}
\textbf{Source problem code:} catalogue entry
\subsection{Definitions and mathematical statement}
An \(A_k\) singularity, for an integer \(k\geq 1\), is a complex plane curve germ analytically equivalent to
\[
\{(u,v)\in\C^2:u^2+v^{k+1}=0\}
\]
at the origin. Equivalently, after an invertible local analytic change of coordinates, a local defining equation is a nowhere-vanishing analytic factor times \(u^2+v^{k+1}\). This specifies the entire local curve germ, rather than merely one branch of a larger germ.

The supplied problem asks for the largest possible \(k\) for a plane curve of degree \(d\), and, more generally, for a curve in \(\A^2(\C)\) of bidegree \((a,b)\), where \((a,b)\in\N^2\).

A precise conditional affine formulation can be given as follows. For a nonzero polynomial \(f\in\C[x,y]\), let \(V(f)\) denote its zero curve with the local defining equations furnished by \(f\). Define
\[
\begin{aligned}
\mathcal K_d^{\mathrm{aff}}
={}&\bigl\{k\in\Z_{\geq1}:\exists f\in\C[x,y]\setminus\{0\},\ \exists p\in\C^2,\\
&\hspace{35mm}\deg f=d,\ f(p)=0,\ (V(f),p)\text{ has type }A_k\bigr\}.
\end{aligned}
\]
Under the convention that bidegree means the exact pair of separate polynomial degrees \((\deg_x f,\deg_y f)\), define
\[
\begin{aligned}
\mathcal K_{a,b}^{\mathrm{bi}}
={}&\bigl\{k\in\Z_{\geq1}:\exists f\in\C[x,y]\setminus\{0\},\ \exists p\in\C^2,\\
&\hspace{12mm}\deg_x f=a,\ \deg_y f=b,\ f(p)=0,\ (V(f),p)\text{ has type }A_k\bigr\}.
\end{aligned}
\]
The requested extremal quantities, under these conventions, are
\[
\max\mathcal K_d^{\mathrm{aff}}
\qquad\text{and}\qquad
\max\mathcal K_{a,b}^{\mathrm{bi}},
\]
whenever the corresponding sets have largest elements. If a set is empty, there is no such singularity and no maximum is assigned here.

These formulas are conditional interpretations, not a resolution of the source's missing conventions. The first question does not explicitly specify the ground field or whether its plane is affine or projective; the complex affine formulation above is suggested by the second question. A complex projective formulation would instead quantify over degree-\(d\) homogeneous equations and points of \(\PP^2(\C)\). The record also does not specify whether curves must be globally reduced or irreducible, whether bidegree is exact or bounded, or whether \(0\in\N\). The displayed polynomial formulation imposes no global reducedness or irreducibility condition, although the prescribed local germ is reduced. These choices require confirmation before a canonical formulation is accepted.
\subsection{Short English statement}
Determine how large an \(A_k\) singularity can occur on a plane curve of degree \(d\), and more generally on a complex affine plane curve of prescribed bidegree \((a,b)\). The precise curve class and degree conventions remain unresolved. The supplied record labels the problem open; this draft does not independently verify that status.
\subsection{Sources}
\begin{itemize}
\item Local registry: \texttt{data/unsolvedmath/problems.json}, permanent UnsolvedMath ID \(30003866\), title ``Maximum Simple Singularities of Plane Curves,'' source problem code \texttt{catalogue entry}, in the Oberwolfach Reports Open Problems collection. No individual proposer is identified in the supplied record; its proposed year is 2018.
\item Supplied report citation: \emph{Subgroups of Cremona Groups} (2019), \emph{Oberwolfach Reports}, pp.~1685--1743, DOI \texttt{10.4171/owr/2018/28}. Supplied link: \url{https://doi.org/10.4171/owr/2018/28}. The background attributes the problem to workshop open-problem discussions documented in this report.
\item Supplied literature reference: Schneider, \emph{Plane curves and \(A_k\)-singularities}, \url{https://arxiv.org/abs/1808.05133}.
\item Literature assessment dated exactly \texttt{2026-08-21}: the supplied curation labels the problem open and states that the cited work provides formulas and tables in small bidegrees, while the requested general extremal function for arbitrary \(a\) and \(b\) is not determined. This is supplied dated curation derived from the OWR material, not an independently verified literature assessment. Its time-sensitive status and cited claims require independent verification.
\end{itemize}
\end{document}
